\begin{center}
{\bf Lời cảm ơn và Lời với bố.}
\end{center}

Kính thưa Ban tổ chức lễ tang.

Kính thưa các cụ, các ông các bà, các bác các cô các chú, anh chị em bạn bè hàng xóm láng giềng gần xa và thân bằng cố hữu của bố chúng tôi.

Kính thưa các quý vị, đại diện các cơ quan đoàn thể.

Từ khi bố chúng tôi lâm bệnh đến tận những phút giây cuối cùng và trong cả buổi tang lễ hôm nay, chúng tôi đã nhận được rất nhiều sự chia sẻ, giúp đỡ của các y bác sỹ, của bạn bè và người thân.

Chúng tôi không biết nói gì hơn là lòng cảm ơn chân thành của cả gia đình tới các ông bà, chú bác, anh chị em và các quý vị, bạn bè của gia đình.

Trong những ngày bố tôi ốm nặng, ông đã được sự quan tâm đặc biệt của tập thể các bác sỹ, y tá, điều dưỡng viên của khoa hồi sức cấp cứu bệnh viện 354 và khoa cấp cứu A9 bệnh viện Bạch Mai. Chúng tôi sẽ không thể nào quên tấm lòng của các y bác sỹ đã hết lòng cứu chữa cho bố tôi.

Chúng tôi cũng xin cảm ơn các cô chú, bạn bè đã gửi những lời động viên sâu sắc trên báo chí, trên mạng xã hội, qua thư điện tử những hồi ức và kỷ niệm về bố.

Chúng tôi chân thành cảm ơn ĐHQG Hà nội, Đại học công nghệ, Viện Hàn lâm khoa học Việt nam, Viện CNTT và các đơn vị trong ban tổ chức tang lễ, nhà tang lễ bệnh Viện trung ương quân đội 108, đã tận tình giúp đỡ gia đình tổ chức tang lễ được chu toàn.


Sau đây tôi xin được nói đôi lời với bố.

Thưa bố,

Lời nói thiêng liêng nhất con xin cảm ơn bố, cuộc đời bố đã để lại cho mẹ, cho chúng con một tinh thần đầy yêu thương và chính trực.

Bố luôn thẳng thắn nhưng cũng rất hài hước và lãng mạn. Mẹ kể, lần gặp đầu tiên, bố tặng mẹ cuốn truyện Ruồi Trâu với một ánh mắt long lanh. Bố làm những bài thơ tặng mẹ đến ông ngoại cũng mê và chép lại. Mỗi niềm vui trăn trở trong công việc, mỗi bài viết tâm huyết với đất nước bố đều chia sẻ với mẹ. Tình yêu của bố sẽ giúp mẹ luôn vững vàng, như sự mạnh mẽ lạc quan mẹ đã luôn có trong những năm qua, khi chăm sóc sức khoẻ cho bố.

Đối với chúng con, niềm say mê của bố đối với khoa học, những trăn trở của bố đối với đất nước, là lời dạy vô giá. Bố nói với con, cuộc sống cần nhất sự trung thực. Tưởng chừng đơn giản nhưng trung thực bao gồm cả dũng cảm, trung thực với chính mình để có những chính kiến độc lập, và trung thực trong cuộc đời để dũng cảm nói lên những ý kiến tâm huyết.

Những lời phát biểu chân thành của bố thật dũng cảm, giản dị nhưng đầy mạnh mẽ về đất nước - như mong ước của bố - những ý kiến của mình sẽ như một giọt nước nhỏ bé hoà vào nhiều triệu giọt nước khác của dân tộc để tạo thành dòng thác “đổi mới" của đất nước.

Giọt nước nhỏ bé đó, như bố lấy tên, là những bài viết « Nghĩ suy cùng đất nước ». Giọt nước đó sẽ không tan đi khi bố chia xa.

Câu thơ bố viết cho con khi con còn nhỏ từ đất nước Đan Mạch của An-đéc-sen sẽ mãi đồng hành cùng con :\\
{\em 
« Và con ơi,\\
Bố muốn tặng con\\
một chú ngựa rơm\\
Trên gập ghềnh đường xa\\ 
con sẽ bước »\\
}

Một lần nữa, xin gửi lời cảm ơn đến tất cả các ông bà, chú bác, anh chị em và các quý vị, bạn bè của gia đình đã đến dự lễ viếng và đã gửi vòng hoa cũng như giúp đỡ chúng tôi tổ chức buổi tiễn đưa bố tôi ngày hôm nay. Trong lúc tang gia bối rối, không tránh khỏi sơ xuất, chúng tôi rất mong có sự thông cảm và lượng thứ.

Phan Dương Hiệu.

